% Local-ring fragment. Insert in Part D after D.1--D.4.
\subsection{Faithfulness and the openness test for regular local rings}

Throughout this subsection $A,B$ are regular local $\Q$-algebras
essentially of finite type.  We use chosen geometric generic points.
Morphisms and ring homomorphisms are over $\Q$.  Write
$E_A=\pi_1(A)$, $E_B=\pi_1(B)$ and $p_A,p_B$ for their arithmetic
maps to $G_\Q$.  To retain the constants when these maps are not
surjective, a relative outer homomorphism means an equivalence
class of pairs
\[
 (\phi,\gamma),\qquad
 \phi:E_A\longrightarrow E_B\text{ continuous},\quad
 \gamma\in G_\Q,\quad
 p_B\phi=\operatorname{Int}(\gamma)p_A,
 \tag{D.5.0}
\]
where
\[
 (\phi,\gamma)\sim
       (\operatorname{Int}(h)\phi,p_B(h)\gamma),\qquad h\in E_B.
\]
This is the explicit base-point transporter in an outer map over
$\operatorname{Spec}\Q$: a path from the image of the source
geometric point to the target point induces $\phi$, and its
arithmetic projection is $\gamma$; changing the path gives exactly
the displayed equivalence.  We write
$\operatorname{Hom}^{\rm out}_{/G_\Q}(E_A,E_B)$ for these classes.
When $p_B$ is surjective, choose $h$ with $p_B(h)=\gamma^{-1}$.
This reduces the pair to one with $\gamma=1$, and the remaining
conjugations have $p_B(h)=1$.  Thus this definition recovers
strict maps over $G_\Q$ modulo target geometric conjugation in
that case.

The transporter also gives the correct map on algebraic constants
without a geometric connectedness assumption.  For number fields
$A=k$, $B=l$ embedded in $\bar\Q$, the condition on the pair is
$\gamma G_k\gamma^{-1}\subset G_l$, and $\phi$ is its forced
inclusion.  This condition is equivalent to $\gamma^{-1}l\subset k$.
The map from its class to $\gamma^{-1}|_l:l\to k$ is a bijection
onto $\operatorname{Hom}_\Q(l,k)$: its fibers are exactly the
left $G_l$-cosets used in the equivalence relation.  This checks
the constants and the direction of the transporter explicitly.
For a general ring its algebraic constant field gives the
corresponding open arithmetic image; every geometric kernel below
is the kernel of the indicated map to $G_\Q$.

\begin{statement}{D.5}{Proved}{Power-separatedness for finitely generated fields}
For every field $F$ finitely generated over $\Q$,
\[
                 \bigcap_{n\ge1}(F^\times)^n=\{1\}.          \tag{D.5.1}
\]
\end{statement}
\begin{proof}
Let $k$ be the algebraic closure of $\Q$ in $F$, a number field,
and choose a normal proper $k$-model of $F$.
Suppose $a=b_n^n$ for some $b_n\in F^\times$ for every $n$.
Every prime-divisor valuation of $a$ is divisible by every $n$
and thus zero; the same valuations of $b_n$ are zero.
Normality and properness imply $a,b_n\in k^\times$: a rational
function with no codimension-one poles is globally regular, and
global regular functions on this proper geometrically integral
model are constants.  The assertion is immediate at this step
as well when $F=k$.

At every finite prime of $k$, the equality $a=b_n^n$ again forces
the valuations of $a$ and all $b_n$ to vanish.  Thus these elements
are algebraic-integer units.  For every embedding $\sigma:k\to\C$,
\[
 |\sigma(b_n)|=|\sigma(a)|^{1/n}
       \le\max\{1,|\sigma(a)|\}.
\]
There are only finitely many algebraic integers of bounded degree
and bounded conjugates: the coefficients of their monic minimal
polynomials are bounded integers by the elementary symmetric
formulas.  Some $b_n=b_m$ therefore occurs for distinct $n,m$,
so this element, and consequently $a$, is a root of unity.
Let $w$ be the order of the finite group $\mu(k)$ and use
$a=b_w^w$.  The element $b_w$ is also a root of unity and
$b_w^w=1$.  Hence $a=1$.
\end{proof}

\begin{statement}{D.6}{Proved}{The arithmetic fundamental-group functor is faithful on these spectra}
Two morphisms $\operatorname{Spec}A\to\operatorname{Spec}B$
inducing the same relative outer homomorphism (D.5.0) are equal.
No dominance assumption is required for this assertion.
\end{statement}
\begin{proof}
Let the corresponding ring maps be $f,g:B\to A$.  For every
$n\ge1$, the coefficient module $M=\mu_n$ comes from $G_\Q$.
For a pair $(\phi,\gamma)$ the exact coefficient transport is
$M$ with action $p_B\phi\to M$ with action $p_A$,
$m\mapsto\gamma^{-1}m$.  In additive cocycle notation its
pullback is therefore
\[
                  c\longmapsto
       \bigl[a\longmapsto\gamma^{-1}c(\phi(a))\bigr].
 \tag{D.6.1}
\]
The equality $p_B\phi=\operatorname{Int}(\gamma)p_A$ verifies
that this is a cocycle with the required source action.  Replacing
the pair by $(\operatorname{Int}(h)\phi,p_B(h)\gamma)$ changes
(D.6.1) by the coboundary of
$-\gamma^{-1}p_B(h)^{-1}c(h)$, as follows directly from
$c(huh^{-1})=p_B(h)c(u)+(1-p_B(huh^{-1}))c(h)$.
Thus the induced cohomology pullback depends only on (D.5.0).
For an actual scheme morphism this is exactly its pullback of
the sheaf $\mu_n$: its base-point path acts on that sheaf by
the same $\gamma$.  Equality of the two relative outer maps
therefore gives equal pullbacks on $H^1(B,\mu_n)$.
The degree-one \'etale/fundamental-group
comparison and the Kummer sequence give, for every $u\in B^\times$,
\[
                  f(u)/g(u)\in(A^\times)^n
                  \qquad\text{for every }n.
\]
Apply D.5 to $\operatorname{Frac}(A)$ to obtain $f(u)=g(u)$.
For arbitrary $b\in B$, choose $c\in\Q$ such that $b-c$ is a
unit.  Such a $c$ exists because the residue of $b$ excludes at
most one rational number.  The maps agree on that unit and on
$c$, so they agree on $b$.  Thus $f=g$.
\end{proof}

\begin{statement}{D.7}{Proved}{An actual morphism has open image exactly when it is dominant}
For a morphism $f:\operatorname{Spec}A\to\operatorname{Spec}B$,
the image of $f_*:\pi_1(A)\to\pi_1(B)$ is open if and only if
the ring map $B\to A$ is injective.  Consequently there is an
injective map
\[
 \operatorname{Hom}^{\mathrm{dom}}_\Q
       (\operatorname{Spec}A,\operatorname{Spec}B)
 \lhook\joinrel\longrightarrow
 \operatorname{Hom}^{\mathrm{out,open}}_{/G_\Q}
       (\pi_1(A),\pi_1(B)).                       \tag{D.7.1}
\]
\end{statement}
\begin{proof}
Suppose first that $B\to A$ is injective, and set
$F=\operatorname{Frac}B\subset E=\operatorname{Frac}A$.
The algebraic closure $F'$ of $F$ in $E$ is finite over $F$.
With compatible separable closures, restriction $G_E\to G_F$
has image $G_{F'}$: it is closed by compactness and its fixed
field inside $F^{\mathrm{sep}}$ is $E\cap F^{\mathrm{sep}}=F'$.
It is therefore open.  Since $A,B$ are normal, the generic-point
maps $G_E\twoheadrightarrow\pi_1(A)$ and
$G_F\twoheadrightarrow\pi_1(B)$ are surjective.  The commutative
square of these four groups identifies the image of $f_*$ with
the image of $G_{F'}$ under the second surjection.  Surjective
homomorphisms of profinite groups are open, so $f_*$ has open image.

Conversely take $0\ne b\in\ker(B\to A)$.  Such an element lies
in the maximal ideal of $B$, since a unit cannot map to zero.
The element $u=1+b$ is a unit, and its image in $A$ is $1$.
It is nonconstant over $\Q$: a nonzero element algebraic over
$\Q$ has inverse given by its minimal polynomial and is a unit
in every $\Q$-algebra containing it; thus $b$ cannot be algebraic.

Let $k$ be the algebraic closure of $\Q$ in $F=\operatorname{Frac}B$.
Normality gives $k\subset B$, since elements of $k$ are integral
over $B$.  The chosen geometric component has field
$F\bar\Q$, regular over $\bar\Q$.  On a normal proper
$\bar\Q$-model, the nonconstant rational function $u$ has a
prime-divisor valuation $d\ne0$.  For every $n$, adjoining
$u^{1/n}$ has ramification index at least $n/\gcd(n,d)$ there.
But $u$ is a unit on the geometric local scheme, so this Kummer
cover is finite \'etale on that scheme.  Choosing compatible roots
gives a homomorphism
\[
 \theta:\Pi_B\longrightarrow\widehat\Z(1),\qquad
 \Pi_B=\ker(\pi_1(B)\to G_\Q),
\]
with infinite image: the orders of its finite Kummer images have
the displayed unbounded lower bounds.

Write $H=\operatorname{im}f_*$.  The pullback of $u$ is $1$, so
its Kummer homomorphism vanishes on $f_*(\Pi_A)$.
Moreover $H\cap\Pi_B=f_*(\Pi_A)$: the arithmetic compatibility
is $p_Bf_*=\operatorname{Int}(\gamma)p_A$, so any preimage of
an element with trivial arithmetic projection also has trivial
arithmetic projection.
Thus $H\cap\Pi_B\subset\ker\theta$.  If $H$ were open, this
would be an open subgroup of $\Pi_B$, forcing the image of
$\theta$ to be finite, a contradiction.  This proves the converse.
D.6 now gives injectivity of (D.7.1).
\end{proof}

\begin{statement}{D.8}{Open}{The exact fullness assertion for D5}
For these regular local rings, it remains to construct a dominant
morphism from an arbitrary relative outer homomorphism
$(\phi,\gamma)$ of (D.5.0) whose continuous homomorphism
$\phi:\pi_1(A)\to\pi_1(B)$ has open image.
D.5--D.7 prove faithfulness and the exact openness
test for actual morphisms; they do not prove surjectivity of
(D.7.1).
\end{statement}

This scope is the spectra of regular local rings essentially of
finite type specified in cote 149, scans 49--50.  For comparison,
a smooth ample hyperplane section of a proper hyperbolic polycurve
of dimension at least two induces a surjection on fundamental
groups by Lefschetz, while its inclusion is not dominant.  This
proper-variety example is given in
\cite[Example 3.13]{BIR_Carlson2026}.  It does not give a map
between the local spectra in (D.7.1).  The local-ring calculation
exhibits the extra finite \'etale covers explicitly: if a local
ring map kills $b\ne0$, all the covers $z^n=1+b$ become trivial
after pullback and their geometric images have unbounded order.
